\chapter{并行架构与求解器设计}
\label{ch:architecture}

\begin{keybox}
\textbf{本章不是给 OpenMP、MPI 或 GPU 排名，而是建立选型过程。}输入是问题规模、网格结构、求解频率、机器拓扑和维护约束；输出才是 rank/thread/GPU 配置、solver hierarchy 与底层求解策略。\[
\boxed{
\text{问题画像}\rightarrow\text{内存可行性}\rightarrow\text{局部工作粒度}
}
\]
\[
\boxed{
\text{通信/同步拓扑}\rightarrow\text{候选执行模型}\rightarrow\text{benchmark 验证}
}
\]
前面各执行模型章节提供机制，本章只负责把机制组合成决策流程；最终选择必须用\chapref{ch:benchmarking}的公平基准和\chapref{ch:diagnosis}的诊断流程验证。
\end{keybox}

\section{应用阶段图与并行分解}

很多求解器设计失败并不是因为某个 kernel 慢，而是因为\textbf{整个应用的阶段并行结构不一致}。例如一个 particle--field 程序可能同时包含
\begin{equation}
\text{particle transport}
\rightarrow
\text{charge deposition}
\rightarrow
\text{field solve}
\rightarrow
\text{particle push}
\rightarrow
\text{geometry update}.
\end{equation}

field solve 很适合 domain-decomposed MPI，不代表 particle transport 也适合相同 rank decomposition。反过来，粒子阶段可能更适合一个共享场副本加大量线程，而 Poisson solve 又可能希望把空间域分给多个 ranks。

因此在 solver 选型以前，先为每个 phase 写出：

\begin{equation}
\boxed{
\{\text{read state},\ \text{write state},\ \text{natural task},\ \text{preferred ownership},\ \text{required synchronization}\}.
}
\label{eq:architecture-phase-contract}
\end{equation}

如果两个相邻 phase 的 preferred ownership 不同，真正的架构成本不是各自单独的 kernel time，而是 ownership conversion：
\begin{equation}
T_{\mathrm{application}}
=
\sum_q T_q^{\mathrm{phase}}
+
\sum_q T_q^{\mathrm{redistribute/synchronize}}.
\label{eq:architecture-phase-total}
\end{equation}

这与\chapref{ch:distributed-gpu}中 coarse redistribution 的逻辑相同，只是现在发生在不同物理/数值 phase 之间。

\section{数据所有权与执行资源}

一个常见错误是把“rank 数”直接当成“数据怎么存”。更稳定的软件架构应显式区分：

\begin{enumerate}
\item \textbf{application ownership}：主状态由谁持有，哪些对象是权威版本；
\item \textbf{solver ownership}：进入求解器以后，网格/矩阵怎样分区；
\item \textbf{execution placement}：当前 kernel 由哪个 CPU thread/GPU 执行；
\item \textbf{active resource set}：当前 phase/level 实际参与计算的 ranks/threads/devices。
\end{enumerate}

前面源码已经给出真实例子。AMReX 的 \texttt{DistributionMapping} 表示 coarse-grid ownership，而 \texttt{BottomCommunicator()} 可以让真正 bottom solve 只使用一个子 communicator；PETSc GAMG 的 \texttt{createlevel()} 可以在构造 coarse matrix 时减少 active processes。它们都说明：
\begin{equation}
\boxed{
\text{应用的全局 rank 集}
\neq
\text{某个 solver level 必须使用的 active rank 集}.
}
\label{eq:architecture-global-vs-active}
\end{equation}

因此“程序由 64 个 MPI ranks 启动”并不意味着每一级 multigrid 都应该让 64 个 ranks 工作。

\section{问题画像与架构筛选}

\subsection{问题特征描述}
在选择框架前，至少记录四组信息。

第一组是\textbf{数值结构}：全局 unknown 数、每层 box/patch 数、是否 AMR、系数是否强跳变、smoother/coarse strategy。

第二组是\textbf{生命周期}：同一 hierarchy 会重复求解多少次，矩阵/系数多久变化一次，setup 能否复用，solver object 是否长期驻留。

第三组是\textbf{应用耦合}：求解前后的 phase 分别怎样拥有数据，是否需要 particle/mesh/field 之间的重分布，哪些状态允许复制。

第四组才是\textbf{机器结构}：socket、NUMA、GPU、NIC、节点内存和节点间网络。

没有这四组画像，“应该用 MPI 还是 OpenMP”通常没有稳定答案。


\subsection{架构与算法匹配}

给定问题，先回答：
\begin{equation}
\{N,\text{geometry},\text{coefficient},\text{AMR},\text{solve frequency},\text{hardware}\}.
\end{equation}
然后再决定并行模型。多重网格的关键不是把所有框架都叠上去，而是减少不必要的层级。

\paragraph{一个不做“架构排名”的选型例子。}
设某规则三维 Poisson 问题需要反复求解同一 hierarchy，单节点内存足够容纳全部数据。这个信息只能把候选范围缩小到“单节点 CPU shared memory”与“单 GPU/多 GPU 中能够完整驻留数据的方案”，还不能直接选出赢家。接下来分别检查：CPU 端是否已接近内存带宽上限；GPU 端最粗层是否出现严重 small-kernel；setup 是否可复用；数据是否需要频繁往返 host/device。最后用\chapref{ch:benchmarking}的相同数学路径做规模扫描，再用\chapref{ch:diagnosis}解释 crossover。这个例子体现本章的原则：\textbf{架构选择是约束逐步筛选，不是先验打分。}

\section{Solver 生命周期}

从\chapref{ch:software}的源码可以看到，成熟库都会把 setup state 和 apply state 分开。应用层也应明确 solver 生命周期：

\begin{equation}
\boxed{
\text{construct}
\rightarrow
\text{build hierarchy/setup}
\rightarrow
\underbrace{\text{apply/solve}\rightarrow\cdots\rightarrow\text{apply/solve}}_{N_{\mathrm{reuse}}}
\rightarrow
\text{rebuild/destroy}.
}
\label{eq:architecture-solver-lifecycle}
\end{equation}

若每次 Poisson solve 都重建 CSR、KSP、AMG hierarchy 或 communication plan，那么 benchmark 中看到的“solver 慢”可能主要是 architecture/lifecycle 问题，而不是迭代算法问题。

因此接口至少应允许下面几类状态长期存在：

\begin{itemize}
\item geometry/hierarchy 与 level metadata；
\item ownership/communication plan；
\item operator coefficients 或可更新的 coefficient view；
\item smoother/bottom-solver setup；
\item persistent workspace/buffer pool；
\item backend execution resources，例如 stream、thread team 或 communicator。
\end{itemize}

真正需要每次 solve 刷新的通常只是 RHS、initial guess、部分 coefficients 和 convergence state。哪些对象可复用必须由数值依赖决定，而不是由“函数写起来方便”决定。

\section{状态复制与分布式状态}

某些耦合应用会遇到一个重要选项：field/mesh 状态是否在多个 workers/ranks 上复制。

若完整场状态内存为 $M_f$，复制到 $P$ 个 rank 的理想内存成本约为
\begin{equation}
M_{\mathrm{rep}}
\approx
P M_f,
\end{equation}
但 particle/local query 可以完全读取本地副本。

分布式存储则近似为
\begin{equation}
M_{\mathrm{dist}}
\approx
M_f
+
M_{\mathrm{ghost}}
+
M_{\mathrm{metadata}},
\end{equation}
但跨分区的 particle/query/field operation 可能需要通信。

因此复制与分布的基本交换是
\begin{equation}
\boxed{
\text{memory replication}
\quad\leftrightarrow\quad
\text{runtime communication / ownership conversion}.
}
\label{eq:architecture-replicate-distribute}
\end{equation}

场规模相对小、读请求极多且更新频率低时，复制可能有工程吸引力；场规模巨大或更新频繁时，分布式 ownership 更自然。这个判断必须放在整个应用的 phase 图中做，而不能只看 Poisson solver 本身。

\section{候选架构}

\subsection{单节点 CPU}

如果问题能够装入单节点内存，优先研究 shared-memory bandwidth、NUMA 和 vectorization；这些机制已经在\chapref{ch:shared-memory}建立。此时 MPI 仍可能有优势，例如现有 domain decomposition 已高度成熟，但应明确\chapref{ch:mpi}所描述的 rank 数据复制、halo 与同步成本是否值得。

一个合理实验矩阵是：
\begin{equation}
(1\ \text{rank}\times C\ \text{threads}),
(2\times C/2),
(4\times C/4),
\ldots,
(C\times1).
\end{equation}
用同一节点比较，找到 memory bandwidth 与 MPI overhead 的平衡点。

\subsection{多节点 CPU}

多节点 CPU 上，分布式内存机制通常不可避免；基础通信模型见\chapref{ch:mpi}。节点内仍可选择纯 MPI 或 MPI+OpenMP。粗网格上减少 ranks 尤其重要：若每个 rank 的 boxes 很少，OpenMP 线程可能没有足够 tile；若 ranks 太多，coarse-grid halo 与 Allreduce 又会过重。

\subsection{单 GPU}

GPU 是否划算不存在脱离硬件和算法的固定 unknown 阈值。\chapref{ch:gpu}已经解释：问题太小时，kernel launch、setup 和数据迁移等固定成本可能占主导；问题足够大以后，规则 stencil 才更有机会把 HBM 带宽和设备并行度发挥出来。这个 crossover 必须通过\chapref{ch:benchmarking}的规模扫描测出来，而不能写成通用常数。无论规模如何，整个 hierarchy 与 bottom solver 是否尽量 device resident 都是首先要检查的条件。

\subsection{多 GPU}

核心设计从“GPU kernel”升级为“GPU+NIC pipeline”。每 GPU 一 rank 是常见起点，但最粗层不应机械保留所有 ranks。GPU-aware MPI、rank consolidation 和 coarse solver 策略往往比进一步手调 fine-grid kernel 更重要。

\subsection{跨厂商 GPU}

若必须同时覆盖 NVIDIA/AMD/Intel，CUDA-only 的维护成本会升高。可以考虑 HIP/SYCL/Kokkos/OpenMP target，或者像 AMReX 那样在高层保持统一 API、底层使用原生后端\cite{amrexGPUdocs2026}。但任何 performance-portability 选择都应配套跨架构 benchmark，而不是把“编译成功”当作完成。

\subsection{纯 MPI、纯共享内存与 MPI+OpenMP}

设节点有 $C$ CPU cores。至少存在三种不同软件结构：
\begin{equation}
\boxed{
C\text{ ranks}\times1\text{ thread},
\qquad
1\text{ rank}\times C\text{ threads},
\qquad
R\text{ ranks}\times(C/R)\text{ threads}.
}
\end{equation}

这三者不仅是启动参数不同，还会改变 ownership、solver state duplication、halo 数量、NUMA placement 与 coarse-level active set。

纯 MPI 的优点：
\begin{itemize}
\item 内存访问局部性和 ownership 清晰；
\item MPI runtime 负责 rank 间并行；
\item 每 rank 单线程，线程安全问题少。
\end{itemize}
缺点：
\begin{itemize}
\item rank 数多，halo messages 与 metadata 增加；
\item coarse level 上活跃 ranks 过多；
\item 每 rank 复制的 solver 数据结构可能增加内存。
\end{itemize}

纯共享内存最适合“单节点内状态天然共享、热点操作可线程化”的情形，它避免 rank 间 halo 和 per-rank solver duplication，但要求 solver 本身真正具有 thread-parallel kernel；如果某关键阶段仍只有单线程，那么把 MPI ranks 全部收掉只会暴露新的串行瓶颈。

MPI+OpenMP 可以减少 rank 数和网络 endpoints，但需要处理 NUMA、thread scaling、barrier 与 MPI thread level。若只有 rank 主线程调用 MPI，\texttt{MPI\_THREAD\_FUNNELED} 一类模式即可表达这种 ownership；若多个 worker 直接发起 MPI，则需要更强的线程支持，并承担相应 runtime 复杂度。

对于 memory-bound stencil，一个 socket 上 32 threads 未必比 8--16 threads 更快；最佳组合必须同时测量 bandwidth saturation、barrier、NUMA 与每 level 的 useful worker 数。

\section{架构约束模型}

候选架构可以先经过四个硬约束过滤：

\paragraph{内存约束.}
\begin{equation}
M_{\mathrm{state}}
+
M_{\mathrm{hierarchy}}
+
M_{\mathrm{workspace}}
+
M_{\mathrm{replication}}
\le
M_{\mathrm{available}}.
\end{equation}

\paragraph{最小有效工作粒度.}
对每个关键 phase/level，需要
\begin{equation}
\frac{W_{\ell,q}}
{P_{\ell,q}^{\mathrm{active}}}
\gtrsim
g_{q,\min},
\end{equation}
否则固定调度/同步/launch 成本会主导。

\paragraph{所有权转换约束.}
若两个相邻 phase 使用不同 decomposition，应显式预算
\begin{equation}
T_{\mathrm{convert}}
\end{equation}
而不是把它藏进某个 solver call。

\paragraph{生命周期摊销约束.}
昂贵 setup 的方案只有在
\begin{equation}
T_{\mathrm{setup}}
+
N_{\mathrm{reuse}}T_{\mathrm{solve}}
\end{equation}
具有优势时才有意义。

通过硬约束以后，剩余候选再进入\chapref{ch:benchmarking}的受控实验与平台最优实验，最后用\chapref{ch:diagnosis}解释差异。

\section{粒子--场耦合的架构案例}

考虑一个抽象 particle--field 应用。若 field solve 使用 $P$ 个空间分区 MPI ranks，而 particle transport 的自然并行方式是大量共享只读 field 的 workers，就可能出现三种架构：

\begin{enumerate}
\item \textbf{全程空间分区 MPI}：field ownership 一致，但粒子跨分区时需要迁移或远端访问；
\item \textbf{单份共享 field + rank 内线程粒子}：粒子简单，但 field solver 必须具备足够强的 shared-memory parallelism；
\item \textbf{field 分布求解 + solve 后复制/重分布}：以额外 memory 和 conversion cost 换粒子阶段更简单的读取模型。
\end{enumerate}

没有哪一个可以仅凭“Poisson MPI 更成熟”或“OpenMP 没通信”直接胜出。真正比较的是
\begin{equation}
T_{\mathrm{step}}
=
T_{\mathrm{particles}}
+
T_{\mathrm{field}}
+
T_{\mathrm{ownership\ conversion}}
+
T_{\mathrm{geometry/other}}.
\label{eq:architecture-coupled-step}
\end{equation}

这类问题说明：\textbf{局部最优 solver architecture 可能不是全应用最优 architecture。}

\section{本章小结}
本章把架构设计从“硬件/API 选型”提升为应用级所有权设计。先画 phase graph，明确每个阶段读写什么状态、最自然的 task 与 ownership；再设计 solver 生命周期、setup 复用与 active resource set；最后才选择 shared memory、MPI、GPU 或 hybrid backend。

前面源码中的 AMReX \texttt{DistributionMapping/BottomCommunicator} 与 PETSc GAMG coarse-process reduction 都说明 active resource set 可以随 hierarchy 改变。更一般地，不同应用 phase 也可以采用不同并行粒度。最终目标不是让某个 solver kernel 单独最快，而是最小化完整 application step 的 time-to-solution 与 ownership conversion 成本。

\section*{习题}

\subsection*{约束与平衡点推导}
\begin{enumerate}[label=\arabic*.]
\item 一个三维问题有 $800^3$ cells，每 cell 核心数据 96 byte。计算单份数据内存；若 V-cycle hierarchy 额外放大 $8/7$、工作区再乘 2.5，估算总内存并判断 80 GiB GPU 是否可容纳。
\item 全局 $512^3$ 网格在 8、64、512 ranks 下立方分解。计算 local cube 边长、每 rank unknowns 和 surface/volume，作为架构候选的局部粒度输入。
\item CPU 路径每 cycle 30 ms、convergence factor 0.12；GPU 路径每 cycle 8 ms，但每 solve 有 40 ms 数据搬运。分别计算做 1、5、50 次重复 solve 的总时间。
\item 两种配置 A/B：A setup=1 s、solve=0.08 s，B setup=0.2 s、solve=0.11 s。求 break-even reuse 次数，并结合内存限制判断何时配置 A 即使更快也不可选。

\item 建立架构选择的约束链：memory capacity、minimum local work、communication cost、setup amortization。把它写成一组必要条件，并说明任何一个不满足都可淘汰候选配置。
\item 推导“重复求解次数”怎样改变昂贵 setup 与快速 solve 方案之间的 break-even 条件。
\item 证明不存在脱离问题规模和机器参数的“MPI 总比 OpenMP 快”或“GPU 总比 CPU 快”的普适排序：构造两个参数场景使排序相反。
\end{enumerate}

\subsection*{配置扫描}
\begin{enumerate}[label=\arabic*.]
\item 写一个配置扫描脚本，对 rank/thread/GPU 数组合运行固定 benchmark，收集内存、cycle time、iteration count 与 setup time，输出 Pareto 候选。
\item 对同一节点比较纯 MPI、MPI+OpenMP、单 GPU（若可用），确保数值路径一致并绘制 time breakdown。
\item 实现一个简单 architecture advisor：除 unknowns、bytes/cell、节点/GPU/网络外，再输入 phase graph、$N_{\mathrm{reuse}}$ 与 ownership-conversion cost；先排除内存/粒度不可行配置，再按完整 application-step benchmark 数据筛选剩余候选。
\end{enumerate}

\subsection*{架构决策}
\begin{enumerate}[label=\arabic*.]
\item 给定一个假想问题画像（规模、AMR/规则网格、重复求解次数、节点/GPU 数），按“内存可行性 $\rightarrow$ 局部粒度 $\rightarrow$ 通信拓扑 $\rightarrow$ benchmark”完成架构决策，不允许用经验口号直接排名。
\item 对粒子--场耦合应用，分别设计“全程 MPI 空间分区”“共享 field + 线程粒子”“field 分布求解后复制/重分布”三种整体架构，用式~\eqref{eq:architecture-coupled-step}比较 idle resources、内存复制、particle migration 与 ownership conversion。
\item 一个 GPU 方案单次 solve 最快，但 setup 和数据迁移使短任务更慢。说明产品环境中应如何依据任务生命周期而不是单 kernel 峰值做选择。
\end{enumerate}
